\documentclass[conference]{IEEEtran}
\IEEEoverridecommandlockouts
% The preceding line is only needed to identify funding in the first footnote. If that is unneeded, please comment it out.
\usepackage{placeins}
\usepackage{dblfloatfix}
\usepackage{float}
\usepackage{cite}
\usepackage{amsmath,amssymb,amsfonts}
\usepackage{algorithmic}
\usepackage{graphicx}
\usepackage{textcomp}
\usepackage{xcolor}
\def\BibTeX{{\rm B\kern-.05em{\sc i\kern-.025em b}\kern-.08em
    T\kern-.1667em\lower.7ex\hbox{E}\kern-.125emX}}
\begin{document}



% Titulo
\title{Avaliação de Desempenho de Infraestrutura de Computação de Borda para Redes de Sensores Agrícolas Utilizando Teoria de Filas\\
% \thanks{Identify applicable funding agency here. If none, delete this.}
}




% Nome dos integrantes
\author{%
\IEEEauthorblockN{1\textsuperscript{st} João Marcos Sousa Rufino Leal}
\IEEEauthorblockA{\textit{Universidade Federal do Piaui} \\
\textit{UFPI}\\
Picos, Piauí \\
jsousarufinoleal@ufpi.edu.br}
}

\maketitle



% Inicio do documento
\begin{abstract}
A agricultura de precisão demanda processamento eficiente de dados gerados por redes de sensores distribuídos em campo. A computação de borda reduz a latência processando dados localmente, mas o dimensionamento adequado é fundamental para garantir qualidade de serviço. Este trabalho avalia o desempenho de uma infraestrutura distribuída de borda utilizando o simulador Java Modelling Tools (JMT) com modelo M/M/c/K e roteamento Join the Shortest Queue (JSQ). Dois cenários são investigados: escalabilidade vertical (variando contêineres de 2 a 10 por nó) e horizontal (variando nós de 2 a 10). As métricas analisadas incluem tempo médio de resposta, throughput, taxa de descarte e utilização. Os resultados demonstram que ambas as estratégias são eficazes sob alta demanda, com configurações de 10 contêineres ou 10 nós mantendo MRT próximo de 1 segundo e taxa de descarte próxima de zero. Os resultados fornecem diretrizes práticas para dimensionamento de infraestruturas de borda em ambientes agrícolas.
\end{abstract}

\begin{IEEEkeywords}
Teoria de Filas, Computação de Borda, Internet das Coisas, Agricultura de Precisão, Avaliação de Desempenho.
\end{IEEEkeywords}

\section{Introdução}

A agricultura de precisão tem se beneficiado do uso de tecnologias de Internet das Coisas (IoT) para monitoramento em campo~\cite{junaidi2025deep}. Sensores coletam continuamente dados como temperatura, umidade do solo e índices pluviométricos, gerando volumes significativos de informações que necessitam de processamento em tempo real. A análise desses dados permite identificar padrões, detectar anomalias e otimizar recursos como água e fertilizantes.

A computação de borda surge como alternativa ao processamento centralizado em nuvem, permitindo o processamento local próximo às áreas de cultivo~\cite{zhang2025edge, gong2025edge}. A adoção de contêineres oferece isolamento, portabilidade e escalabilidade, permitindo que nós executem múltiplos processos em paralelo. Essa arquitetura reduz a latência e a dependência de conectividade com datacenters remotos, fatores críticos em ambientes rurais.

O dimensionamento adequado da infraestrutura de borda é essencial para garantir qualidade de serviço. O número de nós e contêineres impacta métricas como tempo de resposta, throughput e taxa de descarte. Sistemas subdimensionados resultam em perda de dados, enquanto superdimensionados implicam custos desnecessários. A ausência de diretrizes práticas dificulta a tomada de decisão.

A teoria de filas oferece um arcabouço para modelar sistemas com recursos limitados, permitindo avaliação sem implementação física~\cite{santos2021iot, rocha2025enhancing}. O modelo M/M/c/K representa nós computacionais com múltiplos processadores e capacidade limitada~\cite{li2025modeling}. A política JSQ (\textit{Join the Shortest Queue}) permite simular gateways que distribuem requisições de forma balanceada.

Este trabalho avalia uma infraestrutura distribuída para processamento de dados agrícolas utilizando o simulador JMT com modelo M/M/c/K e roteamento JSQ. Dois cenários são investigados: escalabilidade vertical (variando contêineres de 2 a 10) e horizontal (variando nós de 2 a 10). As métricas analisadas incluem tempo de resposta, throughput, taxa de descarte e utilização.

O artigo está organizado da seguinte forma: a Seção II apresenta os trabalhos relacionados; a Seção III descreve a arquitetura; a Seção IV detalha o modelo; a Seção V apresenta os resultados; e a Seção VI conclui o trabalho.


\begin{table*}[t]
\centering
\caption{Trabalhos Relacionados}
\label{tab:related}
\renewcommand{\arraystretch}{1.3}
\begin{tabular}{|p{2.2cm}|p{3.2cm}|p{1.8cm}|p{2.5cm}|p{2.0cm}|p{1.6cm}|c|}
\hline
\textbf{Referência} &
\textbf{Contexto} &
\textbf{Abordagem} &
\textbf{Métricas Analisadas} &
\textbf{Método de Avaliação} &
\textbf{Escalabilidade} &
\textbf{Fila?} \\
\hline
\cite{atalla2023iot} &
Agricultura IoT &
Experimental &
Conectividade, consumo de energia &
Testbed real &
Horizontal &
Não \\
\hline
\cite{benos2023applying} &
Agricultura IoT &
Revisão &
Latência, volume de dados &
Literatura &
Não avaliada &
Não \\
\hline
\cite{zhang2025edge} &
Edge + Agricultura &
Revisão &
Latência, tempo de resposta &
Literatura &
Não avaliada &
Não \\
\hline
\cite{wang2022edge} &
Edge + IoT Agrícola &
Simulação &
Ocupação de servidores, tempo de fila &
Simulador próprio &
Horizontal &
Não \\
\hline
\cite{santos2021iot} &
IoT / Edifícios &
Modelagem &
MRT, Throughput, Drop Rate &
JMT &
Horizontal &
Sim \\
\hline
\cite{rocha2025enhancing} &
Cloud-Edge + Sensores &
Modelagem &
MRT, Utilização, Throughput &
Modelo analítico &
Vertical &
Sim \\
\hline
\cite{mmck2024performance} &
Cloud Computing &
Modelagem &
Tempo de espera, eficiência &
Modelo analítico &
Não avaliada &
Sim \\
\hline
\cite{li2025modeling} &
Edge Computing &
Modelagem &
Atraso, bloqueio, serviço imediato &
Modelo analítico &
Vertical &
Sim \\
\hline
\textbf{Este Trabalho} &
\textbf{Edge + Sensores Agrícolas} &
\textbf{Modelagem} &
\textbf{MRT, Throughput, Drop Rate, Utilização} &
\textbf{JMT} &
\textbf{Vertical e Horizontal} &
\textbf{Sim} \\
\hline
\end{tabular}
\end{table*}

\section{Trabalhos Relacionados}

Esta seção apresenta oito trabalhos relacionados que utilizam
modelos analíticos ou de simulação para avaliar o desempenho
de sistemas IoT, computação de borda ou redes de sensores. A
Tabela~\ref{tab:related} compara suas principais
características segundo quatro critérios: Contexto de
Aplicação, Modelo de Fila, Ferramenta de Avaliação e Métricas.


No contexto de aplicação, \cite{atalla2023iot} propõe um
ecossistema IoT para gerenciamento de culturas, avaliando a
escalabilidade de redes de sensores, enquanto
\cite{benos2023applying} revisa o uso de sensores IoT e Big
Data na agricultura, evidenciando requisitos críticos de
latência e capacidade de processamento.

Quanto à infraestrutura de borda, \cite{zhang2025edge}
apresenta uma revisão do uso de edge computing em agricultura
inteligente, destacando que o processamento local reduz a
dependência de conectividade com a nuvem. \cite{wang2022edge}
investiga a integração entre IoT e computação de borda,
mostrando que a alocação de recursos impacta métricas como
ocupação dos servidores e tempo de espera.

Em relação ao modelo de fila e ferramenta de avaliação,
\cite{santos2021iot} utiliza o JMT para modelar redes IoT com
modelos M/M/c/K, avaliando MRT, throughput e taxa de descarte.
\cite{rocha2025enhancing} aplica o modelo M/M/c/K para
quantificar o desempenho de sistemas Cloud-Edge, enquanto
\cite{mmck2024performance} avalia modelos M/M/1 em ambientes
de nuvem, fornecendo base teórica para a notação de Kendall.
Por fim, \cite{li2025modeling} propõe uma arquitetura de edge
computing heterogênea, derivando métricas como atraso médio e
probabilidade de bloqueio, demonstrando que o aumento de
núcleos de CPU reduz o tempo de espera.

Diferentemente dos estudos revisados, este trabalho avalia uma
infraestrutura distribuída para redes de sensores agrícolas
utilizando o JMT com modelo M/M/c/K e roteamento JSQ. Dois
cenários são investigados: variação de nós em paralelo e
variação de contêineres por nó, analisando MRT, throughput,
taxa de descarte e utilização.



% ============================================================
% SEÇÃO: ARQUITETURA (versão intermediária)
% ============================================================

\section{Arquitetura}

A arquitetura proposta modela um sistema de processamento
distribuído voltado ao tratamento de dados gerados por redes
de sensores em ambientes agrícolas, seguindo o paradigma de computação colaborativa entre camadas~\cite{yu2025cloud}. O sistema é organizado
em três camadas funcionais: a camada de sensoriamento,
responsável pela coleta de dados em campo; a camada de
roteamento, responsável pela distribuição das requisições
entre os nós de processamento; e a camada de processamento,
onde os dados são efetivamente tratados pelos contêineres
embarcados nos nós computacionais. A Figura~\ref{fig:arquitetura}
ilustra a arquitetura do sistema proposto.

\begin{figure}[H]
    \centering
    \includegraphics[width=\columnwidth]{imagens/Arquitetura.png}
    \caption{Arquitetura do sistema de processamento distribuído
    para redes de sensores agrícolas.}
    \label{fig:arquitetura}
\end{figure}

A camada de sensoriamento é composta por dispositivos instalados
em campo que coletam continuamente dados como temperatura,
umidade do solo, índices pluviométricos e dados de irrigação,
gerando requisições modeladas como um processo de Poisson com
taxa de chegada $\lambda$, refletindo o caráter aleatório e
descentralizado da geração de dados em ambientes agrícolas.

A camada de roteamento é representada por um gateway central
que distribui as requisições entre os nós computacionais
utilizando a estratégia JSQ (\textit{Join the Shortest Queue}),
direcionando cada requisição ao nó com menor fila de espera,
minimizando o tempo médio de resposta e evitando sobrecarga
assimétrica entre os nós.

A camada de processamento é composta por nós de borda, como
Raspberry Pi, equipados com contêineres paralelos para análise
local, consolidação de leituras e detecção de anomalias. Cada
nó é modelado como uma fila M/M/c/K~\cite{rocha2025enhancing},
onde $c$ é o número de contêineres e $K$ a capacidade máxima.
Requisições que chegam com a fila cheia são descartadas
(\textit{drop}).

\subsection{Mapeamento entre Arquitetura e Modelo}

A Tabela~\ref{tab:mapeamento} apresenta o mapeamento entre
os componentes da arquitetura e os elementos do modelo no JMT.

\begin{table}[htbp]
\centering
\caption{Mapeamento entre Arquitetura e Modelo de Simulação}
\label{tab:mapeamento}
\renewcommand{\arraystretch}{1.3}
\begin{tabular}{|p{2.8cm}|p{4.6cm}|}
\hline
\textbf{Componente do Modelo} & \textbf{Elemento da Arquitetura} \\
\hline
Source & Sensores agrícolas gerando requisições \\
\hline
Router (JSQ) & Gateway central de roteamento inteligente \\
\hline
Queue (Máquina) & Nó computacional de borda (Raspberry Pi) \\
\hline
Servidores da fila & Contêineres de processamento do nó \\
\hline
Capacidade $K$ & Buffer máximo de requisições do nó \\
\hline
Sink & Dados processados e entregues ao sistema \\
\hline
\end{tabular}
\end{table}

% ============================================================
% SEÇÃO: MODELO (versão intermediária)
% ============================================================

\section{Modelo}

O modelo analítico proposto é baseado na teoria de filas,
utilizando a notação de Kendall para caracterizar formalmente
o comportamento de cada nó computacional da arquitetura
proposta. A Figura~\ref{fig:modelo} ilustra o modelo de
simulação implementado no JMT.

\begin{figure}[H]
    \centering
    \includegraphics[width=\columnwidth]{imagens/MODELO.png}
    \caption{Modelo de simulação implementado no JMT.}
    \label{fig:modelo}
\end{figure}

\subsection{Notação de Kendall}

A notação de Kendall é amplamente utilizada para descrever
sistemas de filas por meio de uma expressão compacta no
formato $A/S/c/K$~\cite{mmck2024performance}, onde $A$ é a
distribuição do processo de chegada, $S$ a distribuição do
tempo de serviço, $c$ o número de servidores e $K$ a
capacidade máxima do sistema, incluindo requisições em
atendimento e em espera. Neste trabalho, cada nó computacional
agrícola é modelado como \textbf{M/M/c/K}: as chegadas seguem
um processo de Poisson com taxa $\lambda$, refletindo o caráter
aleatório da geração de dados pelos sensores; os tempos de
serviço seguem distribuição exponencial com taxa $\mu$ por
contêiner; $c$ contêineres operam em paralelo; e requisições
que chegam com a fila cheia são descartadas (\textit{drop}),
modelando as limitações de memória de dispositivos de borda
como Raspberry Pi e microcontroladores industriais.

\subsection{Parâmetros do Modelo}

A Tabela~\ref{tab:parametros} apresenta os parâmetros
utilizados na simulação e suas justificativas no contexto
agrícola.

\begin{table}[htbp]
\centering
\caption{Parâmetros do Modelo de Simulação}
\label{tab:parametros}
\renewcommand{\arraystretch}{1.3}
\begin{tabular}{|c|c|p{3.5cm}|}
\hline
\textbf{Parâmetro} & \textbf{Valor} & \textbf{Descrição} \\
\hline
$\lambda$ & 1 a 20 req/s & Taxa de chegada dos dados dos sensores \\
\hline
$\mu$ & 1,25 req/s & Taxa de serviço por contêiner \\
\hline
$c$ & 2, 4, 6, 8, 10 & Número de contêineres por nó \\
\hline
$K$ & 30 & Capacidade máxima da fila por nó \\
\hline
$N$ & 2, 4, 6, 8, 10 & Número de nós em paralelo \\
\hline
Roteamento & JSQ & Join the Shortest Queue \\
\hline
Disciplina & FCFS & First Come, First Served \\
\hline
\end{tabular}
\end{table}

\subsection{Fator de Utilização}

O fator de utilização $\rho$ de cada nó computacional
representa a proporção do tempo em que os contêineres estão
ocupados processando requisições. Para o modelo
M/M/c/K~\cite{li2025modeling}:

\begin{equation}
    \rho = \frac{\lambda}{c \cdot \mu}
    \label{eq:utilizacao}
\end{equation}

\noindent onde $\lambda$ é a taxa de chegada de requisições
ao nó, $c$ é o número de contêineres e $\mu$ é a taxa de
serviço por contêiner. O sistema opera de forma estável
quando $\rho < 1$, condição necessária para que a fila não
cresça indefinidamente. Quando $\rho \geq 1$, o sistema
entra em regime de saturação, com crescimento da fila e
aumento do descarte de requisições.

\subsection{Métricas de Desempenho}

O Tempo Médio de Resposta (MRT) representa o tempo médio
que uma requisição leva desde sua chegada ao sistema até a
conclusão do processamento, incluindo o tempo de espera na
fila e o tempo de serviço. Formalmente, pelo Teorema de
Little~\cite{santos2021iot}:

\begin{equation}
    MRT = \frac{N_{q}}{\lambda_{ef}}
    \label{eq:mrt}
\end{equation}

\noindent onde $N_q$ é o número médio de requisições no
sistema e $\lambda_{ef}$ é a taxa efetiva de chegada,
descontando os descartes. O Throughput ($X$) corresponde
à taxa média de requisições processadas com sucesso por
unidade de tempo, e a Taxa de Descarte indica a taxa média
de requisições perdidas por esgotamento da capacidade da
fila:

\begin{equation}
    X = \lambda \cdot (1 - P_{drop})
    \label{eq:throughput}
\end{equation}

\begin{equation}
    \text{Drop Rate} = \lambda \cdot P_{drop}
    \label{eq:droprate}
\end{equation}

A Utilização ($U$) representa a fração média do tempo em
que os contêineres de um nó estão ativos processando
requisições, conforme a Equação~\ref{eq:utilizacao}.

\subsection{Cenários de Avaliação}

No Cenário 1, o número de contêineres por nó varia entre 2,
4, 6, 8 e 10, com $N = 2$ fixo, avaliando a escalabilidade
vertical. No Cenário 2, o número de nós varia entre 2, 4, 6,
8 e 10, com $c = 2$ fixo, avaliando a escalabilidade
horizontal. Em ambos os cenários, $\lambda$ varia de 1 a
20 req/s, permitindo observar o comportamento do sistema
desde o regime de baixa carga até a saturação completa.
 
O modelo é baseado na teoria de filas com notação de Kendall
para caracterizar o comportamento de cada nó computacional.
A Figura~\ref{fig:modelo} ilustra o modelo implementado no JMT.

\section{Estudo de Caso}

Esta seção apresenta os resultados obtidos por meio das
simulações realizadas no JMT, organizados em dois cenários
de avaliação. Em ambos os cenários, a taxa de chegada $\lambda$
varia de 1 a 20 req/s no eixo X dos gráficos. As métricas
avaliadas são o Tempo Médio de Resposta (MRT), o Throughput,
a Taxa de Descarte (Drop Rate) e a Utilização dos nós
computacionais.

\subsection{Cenário 1}

O primeiro cenário avalia a escalabilidade vertical, variando
o número de contêineres por máquina entre 2, 4, 6, 8 e 10,
com $N = 2$ fixo.

O MRT, apresentado na Figura~\ref{fig:cenario1_mrt}, exibe
transição abrupta entre o regime estável e o saturado. A
configuração com 2 contêineres satura em $\lambda \approx 5$
req/s com MRT atingindo 12 segundos, enquanto a com 4
contêineres satura em $\lambda \approx 10$ req/s. As com 6
e 8 contêineres saturam em $\lambda \approx 15$ e
$\lambda \approx 19$ req/s, respectivamente. A com 10
contêineres mantém o MRT próximo de 1 segundo ao longo de
toda a faixa avaliada, sem atingir saturação.

\begin{figure}[H]
    \centering
    \includegraphics[width=\columnwidth]{Graficos/Unificados 1/tempo_resposta_unificado_1.png}
    \caption{Tempo Médio de Resposta (MRT) para o Cenário 1.}
    \label{fig:cenario1_mrt}
\end{figure}

O throughput, ilustrado na Figura~\ref{fig:cenario1_throughput},
cresce linearmente com $\lambda$ enquanto o sistema opera
abaixo da saturação, com cada configuração atingindo um teto
proporcional à sua capacidade total. A configuração com 2
contêineres satura em aproximadamente 5 req/s, enquanto
configurações com maior número de contêineres continuam
crescendo por mais tempo antes de atingir o teto. Destaca-se
que as configurações com 8 e 10 contêineres apresentam
throughput crescente ao longo de toda a faixa avaliada.

\begin{figure}[H]
    \centering
    \includegraphics[width=\columnwidth]{Graficos/Unificados 1/throughput_unificado_1.png}
    \caption{Throughput para o Cenário 1.}
    \label{fig:cenario1_throughput}
\end{figure}

A taxa de descarte, conforme a Figura~\ref{fig:cenario1_drop},
permanece nula para todas as configurações enquanto o sistema
opera abaixo da saturação, com o início do descarte
coincidindo com o ponto de saturação de cada configuração.
A transição entre zero descarte e alto descarte é abrupta,
refletindo a sensibilidade do sistema à variação de
contêineres.

\begin{figure}[H]
    \centering
    \includegraphics[width=\columnwidth]{Graficos/Unificados 1/drop_rate_unificado_1.png}
    \caption{Taxa de Descarte para o Cenário 1.}
    \label{fig:cenario1_drop}
\end{figure}

As curvas de utilização, mostradas na
Figura~\ref{fig:cenario1_utilizacao}, são progressivamente
mais suaves conforme o número de contêineres aumenta, com
maior separação entre as linhas ao longo de toda a faixa
de $\lambda$. A configuração com 10 contêineres atinge 80\%
de utilização em $\lambda = 20$ req/s, enquanto a
configuração com 2 contêineres já opera a 100\% desde
$\lambda \approx 5$ req/s.

\begin{figure}[H]
    \centering
    \includegraphics[width=\columnwidth]{Graficos/Unificados 1/utilizacao_unificado_1.png}
    \caption{Utilização dos nós para o Cenário 1.}
    \label{fig:cenario1_utilizacao}
\end{figure}

\subsection{Cenário 2}

O segundo cenário avalia o impacto da escalabilidade
horizontal do sistema, variando o número de máquinas em
paralelo entre 2, 4, 6, 8 e 10 nós, mantendo fixo o número
de contêineres por máquina em $c = 2$.

Os resultados de MRT, apresentados na
Figura~\ref{fig:cenario2_mrt}, mostram que a configuração
com 2 máquinas apresenta o maior MRT, atingindo
aproximadamente 12 segundos a partir de $\lambda \approx 5$
req/s, ponto em que o sistema entra em saturação. Conforme
o número de máquinas aumenta, o ponto de saturação se
desloca progressivamente para taxas de chegada mais altas,
evidenciando o benefício da escalabilidade horizontal. A
configuração com 10 máquinas mantém o MRT próximo de 1
segundo ao longo de toda a faixa de $\lambda$ avaliada,
demonstrando que o aumento de nós em paralelo é eficaz para
manter baixos tempos de resposta mesmo sob alta demanda.

\begin{figure}[H]
    \centering
    \includegraphics[width=\columnwidth]{Graficos/Unificados 2/tempo_resposta_unificado_2.png}
    \caption{Tempo Médio de Resposta (MRT) para o Cenário 2.}
    \label{fig:cenario2_mrt}
\end{figure}

O throughput, conforme a Figura~\ref{fig:cenario2_throughput},
cresce linearmente com $\lambda$ enquanto o sistema opera
abaixo da saturação. Após o ponto de saturação, cada
configuração apresenta um teto de throughput proporcional
à sua capacidade total de processamento. A configuração com
2 máquinas satura em aproximadamente 5 req/s, enquanto a
configuração com 4 máquinas satura em torno de 10 req/s.
As configurações com 8 e 10 máquinas continuam crescendo
ao longo de toda a faixa avaliada, indicando que possuem
capacidade suficiente para atender à demanda máxima simulada.

\begin{figure}[H]
    \centering
    \includegraphics[width=\columnwidth]{Graficos/Unificados 2/throughput_unificado_2.png}
    \caption{Throughput para o Cenário 2.}
    \label{fig:cenario2_throughput}
\end{figure}

A taxa de descarte, ilustrada na
Figura~\ref{fig:cenario2_drop}, permanece nula para todas
as configurações enquanto o sistema opera abaixo da
saturação. A configuração com 2 máquinas inicia o descarte
a partir de $\lambda \approx 5$ req/s, atingindo
aproximadamente 15 req/s descartados em $\lambda = 20$
req/s. Cada máquina adicional desloca o início do descarte
para valores mais altos de $\lambda$. As configurações com
8 e 10 máquinas apresentam taxa de descarte próxima de zero
ao longo de toda a faixa avaliada, confirmando sua
capacidade de absorver a carga sem perdas significativas.

\begin{figure}[H]
    \centering
    \includegraphics[width=\columnwidth]{Graficos/Unificados 2/drop_rate_unificado_2.png}
    \caption{Taxa de Descarte para o Cenário 2.}
    \label{fig:cenario2_drop}
\end{figure}

A utilização dos nós, apresentada na
Figura~\ref{fig:cenario2_utilizacao}, cresce
proporcionalmente com $\lambda$ e atinge 100\% no ponto de
saturação de cada configuração. A configuração com 2
máquinas satura em $\lambda \approx 5$ req/s, enquanto as
configurações com maior número de máquinas saturam em
valores progressivamente mais altos. A configuração com 10
máquinas atinge apenas 80\% de utilização no limite da
faixa avaliada ($\lambda = 20$ req/s), indicando que ainda
possui capacidade ociosa disponível.

\begin{figure}[H]
    \centering
    \includegraphics[width=\columnwidth]{Graficos/Unificados 2/utilizacao_unificado_2.png}
    \caption{Utilização dos nós para o Cenário 2.}
    \label{fig:cenario2_utilizacao}
\end{figure}


\section{Conclusão}

Este trabalho avaliou o desempenho de uma infraestrutura distribuída de computação de borda para processamento de dados de sensores agrícolas, utilizando o simulador JMT com modelo de filas M/M/c/K e roteamento JSQ.

Os resultados demonstram que tanto a escalabilidade vertical quanto a horizontal são estratégias eficazes para melhorar o desempenho do sistema. Configurações com maior capacidade de processamento deslocam o ponto de saturação para taxas de chegada mais altas, reduzindo o tempo médio de resposta e a taxa de descarte.

No Cenário 1, a configuração com 10 contêineres manteve o MRT próximo de 1 segundo e taxa de descarte próxima de zero. No Cenário 2, a configuração com 10 nós apresentou comportamento similar, atingindo 80\% de utilização em $\lambda = 20$ req/s.

Os resultados fornecem diretrizes práticas para o dimensionamento de infraestruturas de borda em ambientes agrícolas, permitindo que gestores escolham entre escalabilidade vertical ou horizontal conforme as restrições de custo e disponibilidade de recursos.

Como trabalhos futuros, pretende-se investigar cenários híbridos combinando ambas as estratégias de escalabilidade, além de avaliar o impacto de diferentes políticas de roteamento e distribuições de tempo de serviço não exponenciais.



\bibliographystyle{IEEEtran}
\bibliography{referencias}

\end{document}
